\begin{corollary}[Refuted; persistent one-dimensional forms of A.1, A.6, L1 and L2]
\label{A:A2-sharp-curve-consequences}
For the excellent Dedekind domain in
Theorem~\ref{A:A2-persistent-arithmetic-counterexample}, every nonempty
affine open $U$ satisfies
\[
 \sup_{x\in U}\bigl(\cd_2\kappa(x)+\dim\mathcal O_{U,x}\bigr)=1,
 \qquad \cd_2(U)=2.
\]
Every residue field is nonorderable. Thus the failures of A.1 and
both statements of A.6 persist on every nonempty affine open, with
the smallest positive generic cohomological dimension. The upper
curve bound $\cd_p(R)\leq\cd_p(K)+1$ cannot in general be improved
by any nonempty localization, even for excellent regular domains.
\end{corollary}
\begin{proof}
A closed-point residue field is the union of finite fields whose
successive degrees are eventually two, because each retained prime
is eventually inert in every quadratic step. Its absolute Galois
group is a closed subgroup of $\widehat{\mathbf Z}$ with trivial
pro-$2$ subgroup, so its $2$-cohomological dimension is zero.
Its local ring has dimension one. At the generic point the two
terms are respectively one and zero. These exhaust the points of
the Dedekind scheme, proving the displayed supremum. The preceding
theorem gives $\cd_2(U)=2$, and gives the same value for every
nonempty affine open inside $U$; this also computes the supremum
asked for in A.6(ii). The generic field contains the totally
imaginary number field $\mathbf Q(\sqrt{-35})$, and the closed
residue fields have positive characteristic, so none is orderable.

The proved replacement retains the generic lower witness of
Proposition~\ref{A:spread} and the sharp curve upper bound of
Corollary~\ref{A:generic-plus-dimension}. Theorems
\ref{A:Dedekind-Milnor-criterion} and \ref{A:curve-normalization}
give the exact additional condition for eventual equality on
every noetherian integral curve, without excellence or normality:
one localization must make the simultaneous Milnor residue maps
surjective on every finite separable normalization. The constructed
domain fails that condition already for its own Picard group.
\end{proof}
